\section{The five-cube record}\label{S:five}

Table~\ref{tab:S3cl} lists the three orbits of maximal cliques of $G_5$ (distance $\le2$) with the exact value and a maximiser; Table~\ref{tab:S3max} the two orbits of maximal survivable sets with their certificates. The tetrahedron orbit ($80$ cliques of four cells pairwise at distance $2$) has $V=2$ with the uniform certificate: for the representative $\{(1,1,1,1,1),(1,1,1,-1,-1),(1,1,-1,1,-1),(1,1,-1,-1,1)\}$ the pooled cell is $(1,1,0,0,0)$. The square (a two-dimensional face) has $V=3$ with the maximiser listed in Table~\ref{tab:S3cl}, first-step bound $3$ and certified threshold $3$; the product word of the maximiser has period $4$ and dip $3$, and the opposite-corner alternation of the earlier editions is one of its orderings. Every clique of at most three cells of $G_5$ lies in a radius-one ball, so the census changes only at $L=3$: $1{,}616$ survivable sets (the nonempty subsets of the $32$ balls, which overlap) for $L\le2$ and $1{,}696$ for $L\ge3$, against $1{,}776$ pairwise-admissible sets.
\begin{table}[H]\centering\small
\caption{Maximal cliques of $G_5$ by orbit.}\label{tab:S3cl}
\begin{tabular}{rrrll}
\toprule
size & orbit & $V$ & maximiser $y$ & survivable\\
\midrule
6 & 32 & $3$ & $(-1,1,-1,-1,1)$ & yes\\
4 & 80 & $3$ & $(-1,1,0,0,1)$ & yes\\
4 & 80 & $2$ & $(0,1,0,0,1)$ & no\\
\bottomrule
\end{tabular}

\end{table}
\begin{table}[H]\centering\small
\caption{Five-cube maximal survivable orbits with certificates.}\label{tab:S3max}
\begin{tabular}{rrrrrrrrlr}
\toprule
size & orbit & $V$ & $f(S)$ & $L^\ast$ & cap & antichain & lost after & word period / dip & boundary subsets\\
\midrule
6 & 32 & $3$ & 0 & 0 & -- & 1 & -- & 1 / 0 & 0\\
4 & 80 & $3$ & 3 & 3 & 6 & 4 & 0 & 4 / 3 & 0\\
\bottomrule
\end{tabular}

\end{table}
